\chapter{Getting Started}\label{r01:ch:getting-started}

This chapter installs \ps{} and walks through a complete first kernel: the Jacobi update
of \cref{r01:eq:jacobi}. The later chapters explain each step in detail.

\section{Installation}

\ps{} is a pure Python package that requires Python $\geq 3.10$. It is published on PyPI:
\begin{shcode}
pip install "pystencils~=2.0"            # core package
pip install "pystencils[alltrafos]"      # + islpy, py-cpuinfo
\end{shcode}
The core dependencies are \code{sympy~=1.14}, \code{numpy}, \code{appdirs},
\code{joblib}, \code{pyyaml} and \code{fasteners}. A few optional pieces matter in
practice:
\begin{itemize}
  \item \textbf{A C++ compiler.} The just-in-time (JIT) compiler for CPU kernels calls
        \texttt{g++} (the default on Linux) or \texttt{clang++}. It needs the Python and
        NumPy development headers.
  \item \textbf{\texttt{py-cpuinfo}} is needed to detect the SIMD capabilities of x86
        and ARM Linux hosts (\code{Target.auto_cpu()}). Without it, \ps{} warns and falls
        back to the generic CPU target.
  \item \textbf{CuPy} is needed to \emph{run} CUDA or HIP kernels from Python. You
        can \emph{generate} GPU code without it. \textbf{dpctl} plays the same role
        for SYCL.
  \item The extras \texttt{interactive} (matplotlib, Jupyter, \dots), \texttt{doc},
        \texttt{testsuite} and \texttt{dev} pull in tools for notebooks, documentation,
        tests and linting.
\end{itemize}

\paragraph{In this workspace.}
\ps{} is \emph{not} pip-installed here. The sources in \file{submodules/pystencils/src}
are put on the Python path, and the interpreter of the virtual environment
\file{.venv-rvv} (Python~3.14, NumPy~2.5, SymPy~1.14, py-cpuinfo) is used:
\begin{shcode}
cd my-e4-ws
PYTHONPATH=submodules/pystencils/src .venv-rvv/bin/python -c \
    "import pystencils as ps; print(ps.__version__)"     # -> 2.0b1+4.g9a989e16
\end{shcode}
The listing generator of this manual (\file{./build.sh listings}) uses the same setup.

\section{A first kernel}

\paragraph{Step 1 -- import.} By convention \ps{} is imported as \code{ps}:
\pysnippet{ex01_first_kernel.imports}

\paragraph{Step 2 -- declare the fields.}
A \term{field} is a symbolic stand-in for an array. \code{ps.fields} parses a short
description: two fields named \code{src} and \code{dst}, with element type
\code{double} and two spatial dimensions. The fields have no size yet, so the generated
kernel will work for any array shape.
\pysnippet{ex01_first_kernel.fields}

\paragraph{Step 3 -- write the update rule.}
Indexing a field with a tuple of \emph{relative offsets} gives a symbolic neighbour
access: \code{src[1, 0]} is the east neighbour of the current cell, and
\code{dst.center()} is the current cell itself. A \code{ps.Assignment} pairs a
left-hand side with a SymPy right-hand side:
\pysnippet{ex01_first_kernel.assignment}
The assignment is an ordinary SymPy object. Printing it shows the compass names \ps{}
uses for the offsets (E, W, N, S):
\outsnippet{ex01_first_kernel.assignment}

\paragraph{Step 4 -- generate the kernel.}
\code{ps.create_kernel} runs the code generator. It returns a \code{Kernel} object that
holds the generated function. Without a configuration, it targets a generic CPU with
\code{double} precision:
\pysnippet{ex01_first_kernel.create}

\paragraph{Step 5 -- look at the code.}
\code{kernel.get_c_code()} returns the C code. In a Jupyter notebook,
\code{ps.inspect(kernel, show_cpp=True)} shows it with syntax highlighting.
\Cref{r01:lst:first-kernel} shows the result.

\csnippet[nolol=false,caption={The generated Jacobi kernel (\texttt{Target.GenericCPU}).},label=r01:lst:first-kernel]{ex01_first_kernel.code}

Several details of this output return throughout the manual:
\begin{itemize}
  \item Each field becomes a pointer \code{_data_<name>}. Its sizes and strides become
        integer parameters \code{_size_<name>_<i>} and \code{_stride_<name>_<i>}.
        \code{src} is only read, so its pointer is \code{const}. \code{RESTRICT} is a
        macro the JIT defines as \code{__restrict__}.
  \item The loops skip \emph{one ghost layer} on each side (\code{ctr_0 = 1} up to
        \code{_size_dst_0 - 1}). \ps{} inferred this from the largest neighbour offset
        in the update rule.
  \item The constant $1/4$ was moved out of the loops into \code{__c_1_0o4_0}.
  \item Only the shape of \code{dst} is passed: it is the \emph{archetype} field (the
        alphabetically first domain field) whose sizes define the loop bounds. All domain
        fields must have the same shape, and the JIT checks this when the kernel is
        called.
\end{itemize}
\Cref{r01:fig:sweep} shows what this loop nest does on the array.

\begin{figure}[htbp]
  \centering
  % How the generated Jacobi kernel traverses dst (default C layout: ctr_1 is contiguous).
\begin{tikzpicture}[x=0.62cm,y=0.62cm,font=\footnotesize]
  \fill[ghostcell] (0,0) rectangle (8,6);
  \fill[activecell] (1,1) rectangle (7,5);
  \draw[gridline] (0,0) grid (8,6);
  \draw[domainline] (1,1) rectangle (7,5);
  % inner loops (columns) and the jumps between them
  \foreach \i in {1,...,6} {
    \draw[-{Stealth[length=1.5mm]},thick,psblue!80!black] (\i+0.5,1.2) -- (\i+0.5,4.8);
  }
  \foreach \i in {1,...,5} {
    \draw[dashed,psgray] (\i+0.5,4.8) -- (\i+1.5,1.2);
  }
  % ranges
  \draw[decorate,decoration={brace,mirror,amplitude=4pt}] (1,-0.15) -- (7,-0.15)
    node[midway,below=5pt,align=center,celltext]
      {outer loop \texttt{ctr\_0 = 1 \dots{} \_size\_dst\_0 - 2} ($x$)};
  \draw[decorate,decoration={brace,mirror,amplitude=4pt}] (8.15,1) -- (8.15,5)
    node[midway,right=5pt,align=left,celltext]
      {inner loop ($y$)\\\texttt{ctr\_1 = 1 \dots}\\\texttt{\_size\_dst\_1 - 2}\\[2pt]
       (contiguous\\ in memory)};
  \node[celltext,align=center] at (0.5,6.45) {ghost\\[-2pt]layer};
  % the code
  \node[anchor=west,draw=codeframe,fill=codebg,rounded corners=1pt,inner sep=5pt,
        font=\ttfamily\scriptsize,align=left] at (11.5,3)
    {for (ctr\_0 = 1; ctr\_0 < \_size\_dst\_0 - 1; ++ctr\_0)\\
     \ \ for (ctr\_1 = 1; ctr\_1 < \_size\_dst\_1 - 1; ++ctr\_1)\\
     \ \ \ \ dst(ctr\_0, ctr\_1) =\\
     \ \ \ \ \ \ 1/4 * (src(ctr\_0+1, ctr\_1) + \dots);};
\end{tikzpicture}

  \caption[How the first kernel traverses the array]{How the kernel of \cref{r01:lst:first-kernel} visits the cells of \texttt{dst}. The outer loop runs over \texttt{ctr\_0} ($x$) and the inner loop over \texttt{ctr\_1} ($y$), the contiguous coordinate of a C-ordered NumPy array. The ghost layer (grey) is skipped.}
  \label{r01:fig:sweep}
\end{figure}

\paragraph{Step 6 -- compile and run.}
\code{kernel.compile()} wraps the kernel in a small C++ Python extension module,
compiles it with the host compiler, loads it, and returns a callable. Arrays (and
scalar parameters) are passed as \emph{keyword arguments} named after the fields:
\pysnippet{ex01_first_kernel.run}

\paragraph{Step 7 -- check the result.}
For the interior cells, the result matches a NumPy implementation of the same formula.
The ghost layer of \code{dst} remains untouched:
\pysnippet{ex01_first_kernel.check}

\begin{note}[Compilation cache]
  Compiled modules are cached in \file{~/.cache/pystencils/cpujit}, keyed by a SHA-256
  hash of the code and the compiler command line. A second \code{compile()} of the same
  kernel therefore costs only the import of the cached module. \Cref{r01:sec:jit} shows how
  to change this.
\end{note}

\section{What just happened}

Together, these steps are the whole user workflow of \ps:
\begin{center}\small
\begin{tikzpicture}[node distance=0.45cm,
    st/.style={draw=psblue!70,fill=psblue!6,rounded corners=2pt,align=center,
      inner sep=4pt,font=\footnotesize},
    arr/.style={-{Stealth[length=2mm]},thick,draw=psgray}]
  \node[st] (a) {fields};
  \node[st,right=of a] (b) {assignments};
  \node[st,right=of b] (c) {\texttt{create\_kernel}\\+ config};
  \node[st,right=of c] (d) {\texttt{Kernel}};
  \node[st,right=of d] (e) {\texttt{compile()}};
  \node[st,right=of e] (f) {call with\\arrays};
  \foreach \x/\y in {a/b,b/c,c/d,d/e,e/f} \draw[arr] (\x) -- (\y);
\end{tikzpicture}
\end{center}
The following chapters go through them in order:
\Cref{r01:ch:symbolic-frontend} covers everything you can write on the symbolic side,
\cref{r01:ch:kernel-creation} the configuration of the code generator, and
\cref{r01:ch:outputs-execution} the generated artifacts and how to execute them.
